\documentclass[11pt,a4paper]{article}
\usepackage[T1]{fontenc}
\usepackage[utf8]{inputenc}
\usepackage[authoryear,round]{natbib}
\usepackage[english]{babel}
\usepackage[a4paper,margin=2cm]{geometry}
\usepackage{lmodern}
\usepackage{microtype}
\usepackage{array}
\usepackage{booktabs}
\usepackage{xltabular}
\usepackage[hidelinks]{hyperref}

% One paper: citation key, research question, method, result.
\newcommand{\papier}[4]{%
  \citet{#1} & #2 & #3 & #4 \\
  \addlinespace[0.7em]
}

% Reusable table with automatic page breaks and repeated headings.
\newcommand{\tableaurevue}[1]{%
  \begingroup
  \small
  \renewcommand{\arraystretch}{1.2}
  \setlength{\tabcolsep}{5pt}

  \begin{xltabular}{\textwidth}{
    @{}
    >{\raggedright\arraybackslash}p{0.18\textwidth}
    *{3}{>{\raggedright\arraybackslash}X}
    @{}
  }
    \toprule
    \textbf{Citation} &
    \textbf{Research question} &
    \textbf{Method} &
    \textbf{Result} \\
    \midrule
    \endhead

    #1

    \bottomrule
  \end{xltabular}
  \endgroup
}


\title{Literature Review}
\author{Your name}
\date{\today}

\begin{document}
\maketitle

\section{Scope}
Define your research question, inclusion criteria, and the themes
used to organize the review.

\section{Comparison of papers}
The following entries illustrate the format using two energy papers.
Replace them with your own reading notes.

\tableaurevue{
  \papier
    {canalesAssessingTemporalComplementarity2020}
    {How can temporal complementarity among three variable
     renewable energy sources be quantified?}
    {Correlation coefficients and compromise programming;
     application to wind, solar and hydropower in a region of Poland.}
    {The chosen temporal scale strongly affects the estimated
     complementarity index. The application is context-specific.}

  \papier
    {cauzEmbracingWindPower2025}
    {How can wind complement solar PV in the Swiss energy system?}
    {Analysis of Swiss generation data from 2020--2023 using
     Spearman correlations and energy-flow simulations with storage.}
    {Wind and solar show temporal complementarity across several
     time scales. Scenario results depend on the generation mix
     and storage assumptions.}
}

% Add another theme with a new section and \tableaurevue{...}.
% Each row: \papier{BibTeX key}{Question}{Method}{Result}
% Escape literal special characters in cell text: \%, \&, \_.

\section{Synthesis and research gaps}
Compare findings across papers, discuss methodological limitations,
and identify questions that remain open.

\bibliographystyle{plainnat}
\bibliography{references}

\end{document}
